\chapter{Ngân hàng câu hỏi}

Mỗi lần thầy cô duyệt một đề, các câu của đề ấy vào ngân hàng câu hỏi. Thầy cô dùng lại chúng cho đề sau
(Chương~\ref{chuong:soan-de}, mục \emph{Thêm câu từ ngân hàng}).

\viec{Xem ngân hàng}

\anh{T/35-ngan-hang}{}

Bấm \nut{Ngân hàng câu hỏi} ở thanh trái. Mỗi hàng là một câu, kèm mục kiến thức và khối, môn.

\viec{Sửa một câu trong ngân hàng}

\anh[0.45\linewidth]{T/36-sua-ngan-hang}{}

\begin{buoc}
  \item Bấm vào câu cần sửa. Khung \nhan{Sửa câu trong ngân hàng} mở bên phải.
  \item Sửa \nhan{Đề bài}, \nhan{Phương án} hay \nhan{Cách giải}, rồi bấm \nut{Lưu}.
\end{buoc}

\luuy{Chỉ bản trong ngân hàng thay đổi. Đề đã phát hành vẫn giữ câu cũ; đề sau lấy câu này sẽ dùng bản
mới.}
